% Original application; device mechanisms remain in their theory sections.
\subsection{বাস্তব সমস্যা: Lab চলে, Internet চলে না}
\chTwoPair{\textbf{সমস্যা:} ল্যাবের কম্পিউটারগুলো সুইচ দিয়ে স্থানীয় ফাইল ব্যবহার করতে পারছে। কিন্তু ইন্টারনেট সংযোগ পাওয়া যাচ্ছে না।\par
\textbf{সমাধান:} কম্পিউটারের IP কনফিগারেশন, Default Gateway, রাউটার ও ঊর্ধ্বমুখী সংযোগ পরীক্ষা করতে হবে। স্থানীয় সংযোগ সচল থাকায় শুধু সুইচ বদলালে সমাধান নাও হতে পারে।}{\textbf{নতুন PC-র ক্ষেত্রে:} সুইচের MAC Table-এ প্রাপকের ঠিকানা শুরুতে অজানা থাকতে পারে। সুইচের অংশে বর্ণিত ঠিকানা শেখা ও ফরওয়ার্ড করার নিয়ম এখানেও প্রযোজ্য। সার্ভার বা অ্যাপ্লিকেশন অচল কি না, তা পৃথকভাবে পরীক্ষা করতে হবে।}

\textbf{LAN/WAN-এ ডিভাইসের ভূমিকা:} স্থানীয় ডিভাইস সংযুক্ত করতে হাব বা সুইচ এবং নেটওয়ার্কগুলোর মধ্যে সংযোগ দিতে রাউটার বা গেটওয়ে ব্যবহার করা হয়। NIC-ও LAN-এ ডিভাইস যুক্ত করে। তাই প্রশ্নে শুধু LAN Device উল্লেখ থাকলে অন্য সব যন্ত্র বাদ দেওয়া যায় না। MAC ফ্রেম, IP প্যাকেট, প্রোটোকল রূপান্তর ও সংকেতের কাজ দিয়ে পার্থক্য নির্ধারণ করতে হয়।
